\section{Main findings}
\label{sec:conclusions_findings}

The findings below directly address the three research objectives.

Regarding station planning, this thesis developed a flexible station-placement optimization framework and applied to Barcelona (Chapter~\ref{ch:station_planning_in_barcelona}). The \ac{mcdm} scoring component enables the evaluation of candidate sites under heterogeneous planning priorities, allowing different--and potentially conflicting--objectives to be expressed transparently through scenario-specific criteria and weights. This was illustrated through three planning scenarios, each producing station layouts consistent with the intended goals of that scenario. In addition, the framework evaluates configurations using network-level distance metrics computed on a directed, topography-adjusted street graph, providing a more realistic basis for comparison than euclidean or flat-network distances. In practice, slope-adjusted distances change which locations are accessible, leading to relatively greater allocation of stations within steeper areas than would be obtained under flat-distance measures. Finally, the framework supports incremental network growth, as demonstrated in the l'Eixample expansion case study.

Turning to day-to-day operations, the thesis first documents the spatiotemporal mobility patterns of \textit{Bicing} (Chapter~\ref{ch:bicing_mobility_patterns}), characterizing how \acp{e-b} and \acp{m-b} differ in trip counts, duration, distance, speed, and spatial distribution across topography. This mobility documentation provides the context for the remaining contributions: battery-status predictions and \ac{pm}.

The first of these, battery-status forecasting (Chapter~\ref{ch:mobility_and_battery_modeling_bicing}), documents pronounced spatiotemporal variation in station-level battery status associated with neighborhood characteristics. The time between consecutive trips at a station and a distance-based index (combining trip distances and origin frequencies) were found to be moderately positively correlated with station-level average battery (Pearson $r=0.39$ and $r=0.47$, respectively), suggesting that stations where bikes remain docked longer or that receive fewer trips exhibit higher average charge. It then presents a joint forecasting framework that couples a Markov-chain mobility model with a battery model. The battery component combines a physics-inspired consumption model with a linear charging model that achieved low prediction error in validation (RMSE $=0.0178$, approximately 2\%). When coupling the mobility and battery components, the framework supports both bike-level and system-level forecasting. At the system level, validation showed notable short-horizon performance: the Pearson correlation between observed and predicted values reaches $\approx 0.7$ for both station-level availability and average battery at a 30-minute horizon, with correlation decreasing as the forecast horizon increases. The framework therefore provides joint forecasts that can inform charging and rebalancing decisions.

The second, predictive maintenance (Chapter~\ref{ch:failure_forecasting_bicing}), builds on that mobility context to analyse maintenance data. Brake pads, wheel spokes, and chains exhibit different failure dynamics: \acp{e-b} generally undergo more repairs per bike than \acp{m-b} for brakes and spokes, whereas chains show the opposite pattern. Across components, cumulative distance and bike type emerge as the main drivers of wear, confirming distinct maintenance patterns by bike type and component. Building on these empirical findings, the chapter trains \ac{sa} and \ac{ml}-based models to estimate failure risk for the three components. \ac{deepsurv} achieved the best test performance (RMSE 28.75, 18.83, and 43.62 days for brake pads, wheel spokes, and chains, respectively; \mbox{R\textsuperscript{2}} 0.92--0.94). Notably, performance for chains is weaker than for brake pads and spokes, plausibly due to fewer breakdown events and thus less informative training data. Interpretability analyses consistently identify cumulative distance, average speed, and bike type as the most influential predictors, aligning with the empirical failure patterns. Overall, the models are suitable for fleet-wide, component-level risk estimation and can support workshop scheduling and spare-parts planning.